\hwhat{Can every operator input be summarized by one triple: field strength $h$, catalytic strength $\kappa$ and read-out delay $d$? If so, a model fitted on three input classes should predict the fourth. Each model call was coded as idle, work or talk. Every nudge, broadcast human message, named human message and @-mention was timed at the call that first read it (DQ1 ledger). One generalized linear response with a shared kernel and a shared field direction was fitted per period, and each class in turn was held out (5 day folds). Synthetic validation came first. Replication: G51, G38, G04; natives: NE43, \#5 (G05), NE38. \textbf{The claim fails.} Held-out-class skill is 0.47 (nudges) and 0.72 (mentions) of a class-specific fit in G51, 0.06--0.75 in G04, 0.12 for G38 nudges. What holds is narrower: at the read-out itself (0--2 calls) the shared shape keeps 77--100\% of the skill in G51 and for G04 humans. The classes differ in their tails. Holdout not run.}

\hmodel{Models 02 (kinetic response; field vs.\ catalyst split of a rate change, as in H39) and 09 (exogenous kicks with a response kernel, on the call clock). Transition log-odds against staying get a baseline (agent, day, aging, Markov-2) plus a kick term. H59 constrains each class to a uniform catalyst $\kappa_c$ plus a field $h_c$ along one shared direction $\theta$, with one shared kernel $K_\ell$ over lags since the receiving call. Rivals: class-specific free responses (36 numbers per class), a class-specific direction, the read-out delay alone (one amplitude for all) and the inert model.}

\hmath{\vspace{-5pt}\begin{gather*}
\eta_{ij}=\alpha_{ij}\!\cdot\!x+\textstyle\sum_{c,\ell}D_{c\ell}\,\beta_{ij}(c,\ell)\\
\text{H59: }\beta_{ij}(c,\ell)=K_\ell\big[\kappa_c+\tfrac{h_c}{2}(u_j-u_i)\big],\ u=(0,\cos\theta,\sin\theta)\\
T_c=S_{\rm H59}/S_{\rm free},\quad S=\text{held-out log-lik gain over inert}
\end{gather*}\vspace{-7pt}}

\hobs{../figures/summary_obs.pdf}{(a) Held-out skill of each rival relative to the class-specific free fit, per period and class (informative cells only). H59 (orange) is below the 0.8 line everywhere; the read-out-delay rival (blue) is far below for message classes. (b) Post hoc split: at the read-out (rows read 0--2 calls ago) H59 transfers for G51 nudges and mentions and G04 human messages; in the tail (3--30 calls) it does not.}

\htelos{Partly useful. At the moment an agent reads it, an operator input does nearly the same thing whatever its class: the agent is pulled toward talking, with a little extra traffic. A nudge and an @-mention look alike at the receiving call. So one number per class plus the delay until the next model call describes most of the immediate effect. What differs is what follows. Nudged agents stay idle-prone for 30 calls, which looks like the nudger choosing stuck agents rather than a lever effect. Regime-I mentions speed up the work$\leftrightarrow$wait edge. So for planning, treat an input's immediate pull as portable and its aftermath as class-specific. Reading several messages at once counts as one read (G05), and the class levers do not change when the nudger is switched off (NE43).}

\hbottom{One triple per input class does not predict a held-out class. Only the response at the read-out call is shared; the classes differ in their tails and, in some regimes, in which transition they act on.}

\hresults{\footnotesize\begin{tabular}{@{}p{0.37\columnwidth}@{\ }p{0.5\columnwidth}@{\ }c@{}}
Prediction & Observed & \\\hline
\raggedright P1 LOCO $T\ge0.8$ (G51) & \raggedright nudges 0.47, mentions 0.72 & $\times$\\
\raggedright P2 delay alone loses & \raggedright messages yes; nudges tie & $\sim$\\
\raggedright P3 kernel at read-out & \raggedright $K_1$ 0.44 (G51), 0.65 (G38) & $\sim$\\
\raggedright P4 no inbox decay & \raggedright nudges, mentions yes; broadcast humans 0.57 & $\sim$\\
\raggedright P5 lever-table signs & \raggedright all but broadcast-human field size & \checkmark\\
\raggedright P6 G38, G04 & \raggedright both fail & $\times$\\
\raggedright P7 synthetic & \raggedright recovery yes; direction blind at idle & $\sim$\\
\raggedright NE43 invariance & \raggedright $\Delta$ CIs $\ni0$; transfer 1.01 & \checkmark\\
\raggedright G05 dose saturates & \raggedright $\ge$3 vs 1 msg: $-0.02$; LOCO fails & $\sim$\\
\end{tabular}}

\hobsb{../figures/synthetic_validation.pdf}{Synthetic validation with the real exposure design. (a) $\kappa$ and $h$ are recovered at G51 and G38 counts. (b) LOCO verdicts by planted world: a one-lever world passes at G51. A direction violation is caught at G38 but not at G51, because nudges are read almost only by idle agents, where two parameters fit any direction.}

\hcaveats{The read-out/tail split is post hoc (frozen for the holdout as C1). The state space is behavioral only; content fields are not in it. The baseline is fitted once and held fixed in the bootstraps. Named-human (92 reads) and G38 broadcast-human cells are uninformative. Kickoffs are not testable at call level. The nudge tail may be the nudger's targeting, which the past-only baseline cannot see.}

\hnext{Run \texttt{confirm.py} on the \#51 tail. Test whether the nudge tail is selection (trigger state, first vs re-fire). Try a two-component lever (edge catalyst + second direction). Add content drift on the same clock. Pool kickoffs across regime-III periods.}
